% --- Archon physics formalization source begin ---
% archon:physics
% archon:covers PhyXMiniProblems/problem_phyx_mini_0227.lean
% archon:source-report reports/phyx_mini/problem_phyx_mini_0227.source.json

\chapter{Physics problem phyx\_mini\_0227}
\label{ch:PhyXMiniProblems_problem_phyx_mini_0227}

\paragraph{Problem source.}
\#\# Physical scenario

A very light rigid rod of length \textbackslash{}( 0.500\textbackslash{},\textbackslash{}mathrm\{m\} \textbackslash{}) extends straight out from one end of a meterstick. The combination is suspended from a pivot at the upper end of the rod as shown in figure. The combination is then pulled out by a small angle and released.

\#\# Question

By what percentage does the period differ from the period of a simple pendulum \textbackslash{}( 1.00\textbackslash{},\textbackslash{}mathrm\{m\} \textbackslash{}) long?

\#\# Answer choices

- A: A: 2.08\%

- B: B: 3.08\%

- C: C: 5.08\%

- D: D: 4.08\%

\#\# Figure caption (auxiliary; use the image as primary evidence)

The image shows a pendulum setup. A string or rod is attached to a fixed point on a horizontal surface at the top, allowing it to swing freely. Attached to the lower end of the string is a ruler or rod. The pendulum hangs at an angle from its resting vertical position. There is an arrow indicating the length of the pendulum, labeled as "0.500 m." The ruler has measurement markings, indicating that it is used for measuring or indicating length.

\#\# Dataset metadata

Wave Properties; Physical Model Grounding Reasoning, Multi-Formula Reasoning

\paragraph{Recorded answer/context.}
D: D: 4.08\%

\paragraph{Figure/image path.}
/root/proposal\_for\_physic/science-mango/phyx\_mini\_run/phyx\_data/test\_image/227.png

\paragraph{Formalization target.}
create a compiling Lean file with sorry bodies at `PhyXMiniProblems/problem\_phyx\_mini\_0227.lean`. The Lean declarations must preserve the physical quantities, dimensions or dimensional roles, figure labels, governing-law hypotheses, and final relation expressed by this problem.
Use Mathlib/Physlib names found through LeanExplore where available. If a physics API is missing, introduce faithful local abstractions rather than scalar placeholder aliases.

\begin{theorem}[Physics formalization target]
\label{thm:physics:phyx_mini_0227:target}
\lean{PhyXMiniProblems.ProblemPhyXMini0227.periodDifference_matches_recordedAnswerD}
\uses{def:physics:phyx-mini-0227:phyxminiproblems-problemphyxmini0227-timeinseconds, def:physics:phyx-mini-0227:phyxminiproblems-problemphyxmini0227-meterstickpendulumsetup, def:physics:phyx-mini-0227:phyxminiproblems-problemphyxmini0227-matchesmeterstickpendulumfigure, def:physics:phyx-mini-0227:phyxminiproblems-problemphyxmini0227-hasmeterstickpendulumproblemdata, def:physics:phyx-mini-0227:phyxminiproblems-problemphyxmini0227-satisfiespendulumlaws, def:physics:phyx-mini-0227:phyxminiproblems-problemphyxmini0227-perioddifferencepercent, def:physics:phyx-mini-0227:phyxminiproblems-problemphyxmini0227-recordedanswerchoice, def:physics:phyx-mini-0227:phyxminiproblems-problemphyxmini0227-roundstodisplayedpercentage, def:physics:phyx-mini-0227:phyxminiproblems-problemphyxmini0227-isclosestdisplayedpercentage}
The assigned autoformalize agent should translate this physics problem into Lean declarations in the covered file, with theorem and lemma proof bodies written as `by sorry`.
\end{theorem}
\begin{proof}
This is an autoformalization task, not a proof task. Produce statements that can later be proved by the physics prover without weakening the physical model.
\end{proof}

% --- Archon declaration topology begin ---
\section{Declaration topology}
\begin{definition}[Length Quantity]
  \label{def:physics:phyx-mini-0227:phyxminiproblems-problemphyxmini0227-lengthquantity}
  \lean{PhyXMiniProblems.ProblemPhyXMini0227.LengthQuantity}
  A physical length, independent of the unit system used to read it
\end{definition}
\begin{proof}
  This is the declared model definition of Length Quantity.
\end{proof}
\begin{definition}[Time Quantity]
  \label{def:physics:phyx-mini-0227:phyxminiproblems-problemphyxmini0227-timequantity}
  \lean{PhyXMiniProblems.ProblemPhyXMini0227.TimeQuantity}
  A physical time interval
\end{definition}
\begin{proof}
  This is the declared model definition of Time Quantity.
\end{proof}
\begin{definition}[Mass Quantity]
  \label{def:physics:phyx-mini-0227:phyxminiproblems-problemphyxmini0227-massquantity}
  \lean{PhyXMiniProblems.ProblemPhyXMini0227.MassQuantity}
  A physical mass
\end{definition}
\begin{proof}
  This is the declared model definition of Mass Quantity.
\end{proof}
\begin{definition}[Acceleration Quantity]
  \label{def:physics:phyx-mini-0227:phyxminiproblems-problemphyxmini0227-accelerationquantity}
  \lean{PhyXMiniProblems.ProblemPhyXMini0227.AccelerationQuantity}
  A physical acceleration, with dimension length per time squared
\end{definition}
\begin{proof}
  This is the declared model definition of Acceleration Quantity.
\end{proof}
\begin{definition}[Moment Of Inertia Quantity]
  \label{def:physics:phyx-mini-0227:phyxminiproblems-problemphyxmini0227-momentofinertiaquantity}
  \lean{PhyXMiniProblems.ProblemPhyXMini0227.MomentOfInertiaQuantity}
  A scalar moment of inertia about an axis, with dimension mass times length squared
\end{definition}
\begin{proof}
  This is the declared model definition of Moment Of Inertia Quantity.
\end{proof}
\begin{definition}[length In Meters]
  \label{def:physics:phyx-mini-0227:phyxminiproblems-problemphyxmini0227-lengthinmeters}
  \lean{PhyXMiniProblems.ProblemPhyXMini0227.lengthInMeters}
  \uses{def:physics:phyx-mini-0227:phyxminiproblems-problemphyxmini0227-lengthquantity}
  Read a physical length in metres
\end{definition}
\begin{proof}
  This is the declared model definition of length In Meters.
\end{proof}
\begin{definition}[time In Seconds]
  \label{def:physics:phyx-mini-0227:phyxminiproblems-problemphyxmini0227-timeinseconds}
  \lean{PhyXMiniProblems.ProblemPhyXMini0227.timeInSeconds}
  \uses{def:physics:phyx-mini-0227:phyxminiproblems-problemphyxmini0227-timequantity}
  Read a physical time interval in seconds
\end{definition}
\begin{proof}
  This is the declared model definition of time In Seconds.
\end{proof}
\begin{definition}[mass In Kilograms]
  \label{def:physics:phyx-mini-0227:phyxminiproblems-problemphyxmini0227-massinkilograms}
  \lean{PhyXMiniProblems.ProblemPhyXMini0227.massInKilograms}
  \uses{def:physics:phyx-mini-0227:phyxminiproblems-problemphyxmini0227-massquantity}
  Read a physical mass in kilograms
\end{definition}
\begin{proof}
  This is the declared model definition of mass In Kilograms.
\end{proof}
\begin{definition}[acceleration In Meters Per Second Squared]
  \label{def:physics:phyx-mini-0227:phyxminiproblems-problemphyxmini0227-accelerationinmeterspersecondsquared}
  \lean{PhyXMiniProblems.ProblemPhyXMini0227.accelerationInMetersPerSecondSquared}
  \uses{def:physics:phyx-mini-0227:phyxminiproblems-problemphyxmini0227-accelerationquantity}
  Read an acceleration in metres per second squared
\end{definition}
\begin{proof}
  This is the declared model definition of acceleration In Meters Per Second Squared.
\end{proof}
\begin{definition}[moment Of Inertia In Kilogram Meter Squared]
  \label{def:physics:phyx-mini-0227:phyxminiproblems-problemphyxmini0227-momentofinertiainkilogrammetersquared}
  \lean{PhyXMiniProblems.ProblemPhyXMini0227.momentOfInertiaInKilogramMeterSquared}
  \uses{def:physics:phyx-mini-0227:phyxminiproblems-problemphyxmini0227-momentofinertiaquantity}
  Read a scalar moment of inertia in kilogram metre squared
\end{definition}
\begin{proof}
  This is the declared model definition of moment Of Inertia In Kilogram Meter Squared.
\end{proof}
\begin{definition}[Suspension Rod End]
  \label{def:physics:phyx-mini-0227:phyxminiproblems-problemphyxmini0227-suspensionrodend}
  \lean{PhyXMiniProblems.ProblemPhyXMini0227.SuspensionRodEnd}
  The two labeled ends of the very light suspension rod
\end{definition}
\begin{proof}
  This is the declared model definition of Suspension Rod End.
\end{proof}
\begin{definition}[Meterstick End]
  \label{def:physics:phyx-mini-0227:phyxminiproblems-problemphyxmini0227-meterstickend}
  \lean{PhyXMiniProblems.ProblemPhyXMini0227.MeterstickEnd}
  The two distinguished ends of the meterstick
\end{definition}
\begin{proof}
  This is the declared model definition of Meterstick End.
\end{proof}
\begin{definition}[Component Alignment]
  \label{def:physics:phyx-mini-0227:phyxminiproblems-problemphyxmini0227-componentalignment}
  \lean{PhyXMiniProblems.ProblemPhyXMini0227.ComponentAlignment}
  Relative axial arrangement of the support rod and meterstick
\end{definition}
\begin{proof}
  This is the declared model definition of Component Alignment.
\end{proof}
\begin{definition}[Angle Reference Direction]
  \label{def:physics:phyx-mini-0227:phyxminiproblems-problemphyxmini0227-anglereferencedirection}
  \lean{PhyXMiniProblems.ProblemPhyXMini0227.AngleReferenceDirection}
  Reference direction from which the release angle is measured
\end{definition}
\begin{proof}
  This is the declared model definition of Angle Reference Direction.
\end{proof}
\begin{definition}[Stick Mass Distribution]
  \label{def:physics:phyx-mini-0227:phyxminiproblems-problemphyxmini0227-stickmassdistribution}
  \lean{PhyXMiniProblems.ProblemPhyXMini0227.StickMassDistribution}
  The mass-distribution idealization used for the meterstick
\end{definition}
\begin{proof}
  This is the declared model definition of Stick Mass Distribution.
\end{proof}
\begin{definition}[Meterstick Pendulum Setup]
  \label{def:physics:phyx-mini-0227:phyxminiproblems-problemphyxmini0227-meterstickpendulumsetup}
  \lean{PhyXMiniProblems.ProblemPhyXMini0227.MeterstickPendulumSetup}
  \uses{def:physics:phyx-mini-0227:phyxminiproblems-problemphyxmini0227-lengthquantity, def:physics:phyx-mini-0227:phyxminiproblems-problemphyxmini0227-timequantity, def:physics:phyx-mini-0227:phyxminiproblems-problemphyxmini0227-massquantity, def:physics:phyx-mini-0227:phyxminiproblems-problemphyxmini0227-accelerationquantity, def:physics:phyx-mini-0227:phyxminiproblems-problemphyxmini0227-momentofinertiaquantity, def:physics:phyx-mini-0227:phyxminiproblems-problemphyxmini0227-suspensionrodend, def:physics:phyx-mini-0227:phyxminiproblems-problemphyxmini0227-meterstickend, def:physics:phyx-mini-0227:phyxminiproblems-problemphyxmini0227-componentalignment, def:physics:phyx-mini-0227:phyxminiproblems-problemphyxmini0227-anglereferencedirection, def:physics:phyx-mini-0227:phyxminiproblems-problemphyxmini0227-stickmassdistribution}
  All physical quantities and labeled geometric features used in the comparison. The two oscillator fields are scalar models of the linearized angular coordinates. Their Physlib m and k fields are connected to dimensionful generalized inertia and gravitational restoring coefficients only by SatisfiesPendulumLaws; no requested period difference is stored here
\end{definition}
\begin{proof}
  This is the declared model definition of Meterstick Pendulum Setup.
\end{proof}
\begin{definition}[Matches Meterstick Pendulum Figure]
  \label{def:physics:phyx-mini-0227:phyxminiproblems-problemphyxmini0227-matchesmeterstickpendulumfigure}
  \lean{PhyXMiniProblems.ProblemPhyXMini0227.MatchesMeterstickPendulumFigure}
  \uses{def:physics:phyx-mini-0227:phyxminiproblems-problemphyxmini0227-lengthinmeters, def:physics:phyx-mini-0227:phyxminiproblems-problemphyxmini0227-meterstickpendulumsetup}
  Primary-figure readouts and attachment geometry. The image labels the support rod as 0.500 m, places the pivot at its upper end, and shows the meterstick attached at the lower end in the same straight line. Calling the ruled body a meterstick supplies its 1.00 m length; the source does not specify which numbered mark is at the junction, so the end is labeled only as rodSide
\end{definition}
\begin{proof}
  This is the declared model definition of Matches Meterstick Pendulum Figure.
\end{proof}
\begin{definition}[Has Meterstick Pendulum Problem Data]
  \label{def:physics:phyx-mini-0227:phyxminiproblems-problemphyxmini0227-hasmeterstickpendulumproblemdata}
  \lean{PhyXMiniProblems.ProblemPhyXMini0227.HasMeterstickPendulumProblemData}
  \uses{def:physics:phyx-mini-0227:phyxminiproblems-problemphyxmini0227-lengthinmeters, def:physics:phyx-mini-0227:phyxminiproblems-problemphyxmini0227-massinkilograms, def:physics:phyx-mini-0227:phyxminiproblems-problemphyxmini0227-accelerationinmeterspersecondsquared, def:physics:phyx-mini-0227:phyxminiproblems-problemphyxmini0227-meterstickpendulumsetup}
  Textual data and standard idealizations. “Very light” is modeled by zero support-rod mass, while the meterstick is treated as uniform. The proposition releaseIsInSmallAngleRegime retains the source's qualitative small-angle condition without imposing an arbitrary numerical cutoff
\end{definition}
\begin{proof}
  This is the declared model definition of Has Meterstick Pendulum Problem Data.
\end{proof}
\begin{definition}[Satisfies Pendulum Laws]
  \label{def:physics:phyx-mini-0227:phyxminiproblems-problemphyxmini0227-satisfiespendulumlaws}
  \lean{PhyXMiniProblems.ProblemPhyXMini0227.SatisfiesPendulumLaws}
  \uses{def:physics:phyx-mini-0227:phyxminiproblems-problemphyxmini0227-lengthinmeters, def:physics:phyx-mini-0227:phyxminiproblems-problemphyxmini0227-timeinseconds, def:physics:phyx-mini-0227:phyxminiproblems-problemphyxmini0227-massinkilograms, def:physics:phyx-mini-0227:phyxminiproblems-problemphyxmini0227-accelerationinmeterspersecondsquared, def:physics:phyx-mini-0227:phyxminiproblems-problemphyxmini0227-momentofinertiainkilogrammetersquared, def:physics:phyx-mini-0227:phyxminiproblems-problemphyxmini0227-meterstickpendulumsetup}
  The uniform-meterstick, parallel-axis, gravitational-restoring, and small-angle period laws used in the calculation. For the physical pendulum, the generalized inertia is I\_p and the angular restoring coefficient is M g d. For the simple-pendulum comparison, these are m L² and m g L. The reference bob mass is deliberately retained as a physical parameter even though it cancels from the final period ratio. None of these fields states the requested ratio or percentage difference
\end{definition}
\begin{proof}
  This is the declared model definition of Satisfies Pendulum Laws.
\end{proof}
\begin{definition}[period Difference Percent]
  \label{def:physics:phyx-mini-0227:phyxminiproblems-problemphyxmini0227-perioddifferencepercent}
  \lean{PhyXMiniProblems.ProblemPhyXMini0227.periodDifferencePercent}
  \uses{def:physics:phyx-mini-0227:phyxminiproblems-problemphyxmini0227-timequantity, def:physics:phyx-mini-0227:phyxminiproblems-problemphyxmini0227-timeinseconds}
  Absolute percentage by which an observed period differs from a positive reference period. Both arguments are physical times; only coherent SI readouts enter the dimensionless ratio
\end{definition}
\begin{proof}
  This is the declared model definition of period Difference Percent.
\end{proof}
\begin{definition}[Answer Choice]
  \label{def:physics:phyx-mini-0227:phyxminiproblems-problemphyxmini0227-answerchoice}
  \lean{PhyXMiniProblems.ProblemPhyXMini0227.AnswerChoice}
  Labels of the four displayed multiple-choice answers
\end{definition}
\begin{proof}
  This is the declared model definition of Answer Choice.
\end{proof}
\begin{definition}[percent]
  \label{def:physics:phyx-mini-0227:phyxminiproblems-problemphyxmini0227-answerchoice-percent}
  \lean{PhyXMiniProblems.ProblemPhyXMini0227.AnswerChoice.percent}
  \uses{def:physics:phyx-mini-0227:phyxminiproblems-problemphyxmini0227-answerchoice}
  Percentage printed beside each answer label
\end{definition}
\begin{proof}
  This is the declared model definition of percent.
\end{proof}
\begin{definition}[recorded Answer Choice]
  \label{def:physics:phyx-mini-0227:phyxminiproblems-problemphyxmini0227-recordedanswerchoice}
  \lean{PhyXMiniProblems.ProblemPhyXMini0227.recordedAnswerChoice}
  \uses{def:physics:phyx-mini-0227:phyxminiproblems-problemphyxmini0227-answerchoice}
  The answer label recorded in the supplied dataset metadata
\end{definition}
\begin{proof}
  This is the declared model definition of recorded Answer Choice.
\end{proof}
\begin{definition}[Rounds To Displayed Percentage]
  \label{def:physics:phyx-mini-0227:phyxminiproblems-problemphyxmini0227-roundstodisplayedpercentage}
  \lean{PhyXMiniProblems.ProblemPhyXMini0227.RoundsToDisplayedPercentage}
  \uses{def:physics:phyx-mini-0227:phyxminiproblems-problemphyxmini0227-answerchoice}
  A percentage rounds to the same hundredth of a percent as a displayed choice
\end{definition}
\begin{proof}
  This is the declared model definition of Rounds To Displayed Percentage.
\end{proof}
\begin{definition}[Is Closest Displayed Percentage]
  \label{def:physics:phyx-mini-0227:phyxminiproblems-problemphyxmini0227-isclosestdisplayedpercentage}
  \lean{PhyXMiniProblems.ProblemPhyXMini0227.IsClosestDisplayedPercentage}
  \uses{def:physics:phyx-mini-0227:phyxminiproblems-problemphyxmini0227-answerchoice}
  A displayed choice is at least as close as every listed alternative
\end{definition}
\begin{proof}
  This is the declared model definition of Is Closest Displayed Percentage.
\end{proof}
\begin{definition}[Torque quantity]
  \label{def:physics:phyx-mini-0227:phyxminiproblems-problemphyxmini0227-torquequantity}
  \lean{PhyXMiniProblems.ProblemPhyXMini0227.TorqueQuantity}
  A signed physical torque has dimension mass times length squared per time squared.
\end{definition}
\begin{proof}
  This is the dimensionful physical role used by the pendulum laws.
\end{proof}
\begin{definition}[Torque in newton metres]
  \label{def:physics:phyx-mini-0227:phyxminiproblems-problemphyxmini0227-torqueinnewtonmeters}
  \lean{PhyXMiniProblems.ProblemPhyXMini0227.torqueInNewtonMeters}
  \uses{def:physics:phyx-mini-0227:phyxminiproblems-problemphyxmini0227-torquequantity}
  This scalar projection reads a signed torque coherently in SI newton metres.
\end{definition}
\begin{proof}
  Evaluate the dimensionful torque using the coherent SI unit choices.
\end{proof}
% --- Archon declaration topology end ---
% --- Archon physics formalization source end ---
